% !TeX program = xelatex
% ====================================================================
% yuhang_yang_cv.tex -- Bilingual academic CV for Yuhang Yang
% Compile: latexmk -xelatex -gg -interaction=nonstopmode -halt-on-error yuhang_yang_cv.tex
% ====================================================================
\documentclass[11pt]{article}

\usepackage{cv}

\cvname{阳宇航}
\cvheadline{Undergraduate in Information and Computational Science · Nanjing University}
\cvphone{(+86) 193 5846 6910}
\cvlocation{Nanjing, China}
\cvemail{yuhangyang@smail.nju.edu.cn}
\cvwebsite[yuhangyang.site]{https://yuhangyang.site}

\begin{document}

\makecvheader

\cvsection{Education}

\begin{cventry}{2024--2028（预计）}{B.Sc. in Information and Computational Science}{南京大学数学学院}{Nanjing, China}
学位学分绩：\textbf{4.45/5.0}；专业排名：\textbf{1/15}。
\begin{cvitems}
  \item Core Coursework：数学分析（85）、高等代数（92）、解析几何（99）、数值分析（91）、离散数学（95）。
  \item English Proficiency：CET-4 627；CET-6 508。
\end{cvitems}
\end{cventry}

\cvsection{Research Interests}

研究兴趣包括 Optimization Theory、Natural Language Processing（NLP）与 Large Language Models（LLMs），关注数学理论、模型实现与训练行为之间的联系。

\cvsection{Projects}

\begin{cvproject}[https://github.com/YoungDrifter/min-GPT]{2026.08}{min-GPT: A From-Scratch Conversational GPT}{Independent Developer}{PyTorch · GPT-2 · DailyDialog · W\&B}
\begin{cvitems}
  \item 使用 PyTorch 手写 GPT-2 byte-level BPE tokenizer、Causal Self-Attention、MLP 与 Pre-LN Decoder-only Transformer，并加载 DialoGPT-small 预训练权重。
  \item 在 DailyDialog 上采用 reply-only loss 进行 full-parameter fine-tuning；test-set perplexity 由 \textbf{42.34} 降至 \textbf{14.28}。
\end{cvitems}
\end{cvproject}

\cvsection{Competitions}

\cvdateditem{\textbf{Second Prize} · 全国大学生数学竞赛（江苏赛区）}{2025}
\cvdateditem{\textbf{Third Prize} · 全国大学生数学建模竞赛（江苏赛区）}{2025}

\cvsection{Honors}

\cvdateditem{南京大学人民奖学金\textbf{二等奖}（专业前 20\%）}{2025}

\cvsection{Technical Skills}

\cvitem{Languages}{Python、C++}
\cvitem{Frameworks}{PyTorch}
\cvitem{Tools}{\LaTeX、Git、Weights \& Biases（W\&B）}

\end{document}
